\clearpage

\section{Numerical Verification and
Reproduction}\label{sec-nbcp-verification-appendix}

\subsection{Implementation findings and coordinate
checks}\label{sec-legacy}

The archived gap implementation has known coordinate and
Hamiltonian-assembly discrepancies. Their detailed code locations and
interpretation are retained in the
\href{../../../GOVERNMENT/Working-Pad/issue-notes/open/260918-nbcp-physics-code-review.md\#preserved-legacy-implementation-findings}{physics-to-code
review record}. The matrix and coordinate comparison below records the
existing numerical evidence, without treating agreement with the archive
as physical validation.

\subsubsection{Matrix and coordinate
checks}\label{matrix-and-coordinate-checks}

The matrix comparison uses independent Cartesian Holstein-Primakoff (HP)
coefficients at \(J_{PD}=J_\Gamma=0.005\) meV. Matrix errors are in meV;
the two hard-curvature columns are in meV per spin.

\begin{longtable}[]{@{}
  >{\raggedright\arraybackslash}p{(\linewidth - 8\tabcolsep) * \real{0.1579}}
  >{\raggedleft\arraybackslash}p{(\linewidth - 8\tabcolsep) * \real{0.2105}}
  >{\raggedleft\arraybackslash}p{(\linewidth - 8\tabcolsep) * \real{0.2105}}
  >{\raggedleft\arraybackslash}p{(\linewidth - 8\tabcolsep) * \real{0.2105}}
  >{\raggedleft\arraybackslash}p{(\linewidth - 8\tabcolsep) * \real{0.2105}}@{}}
\toprule\noalign{}
\begin{minipage}[b]{\linewidth}\raggedright
Phase
\end{minipage} & \begin{minipage}[b]{\linewidth}\raggedleft
Current H error
\end{minipage} & \begin{minipage}[b]{\linewidth}\raggedleft
Archive H error
\end{minipage} & \begin{minipage}[b]{\linewidth}\raggedleft
Original chart
\end{minipage} & \begin{minipage}[b]{\linewidth}\raggedleft
Equivalent chart
\end{minipage} \\
\midrule\noalign{}
\endfirsthead
\toprule\noalign{}
\begin{minipage}[b]{\linewidth}\raggedright
Phase
\end{minipage} & \begin{minipage}[b]{\linewidth}\raggedleft
Current H error
\end{minipage} & \begin{minipage}[b]{\linewidth}\raggedleft
Archive H error
\end{minipage} & \begin{minipage}[b]{\linewidth}\raggedleft
Original chart
\end{minipage} & \begin{minipage}[b]{\linewidth}\raggedleft
Equivalent chart
\end{minipage} \\
\midrule\noalign{}
\endhead
\bottomrule\noalign{}
\caption{Matrix agreement and hard-coordinate curvature checks against direct HP coefficients.}\label{tab-nbcp-d-numerical-verification-1}\\
\endlastfoot
Y & \(2.784\times10^{-17}\) & \(1.536\times10^{-2}\) & 0.0206658746 &
0.0766217798 \\
V & \(4.164\times10^{-17}\) & \(2.017\times10^{-2}\) & 0.0850517129 &
0.0255063026 \\
\end{longtable}

Changing one spin from (theta, phi) to (-theta, phi+pi) leaves the
physical configuration and the current zero-point energy unchanged, but
reapplying the archived uniform-theta rule selects a different physical
perturbation and hence a different hard curvature. The relaxed
susceptibility is invariant, as checked in the detailed derivation
above.

\clearpage

\subsection{Numerical convergence}\label{sec-convergence}

The reciprocal quadrature uses N x N midpoints and six rotated copies,
without an endpoint duplication or a k=0 regulator. All positive
physical bands and the normal-ordering trace constant enter the vacuum
energy. Fits include all harmonics 1 through 18, without imposing Z3/Z6.
Very small fitted curvature is marked unresolved against an estimated
floating-point/fit floor; this estimate is not a rigorous error bound.

\begin{longtable}[]{@{}rllr@{}}
\toprule\noalign{}
Base mesh N & Phase & SOC = 0.010 meV & Gap (meV) \\
\midrule\noalign{}
\endfirsthead
\toprule\noalign{}
Base mesh N & Phase & SOC = 0.010 meV & Gap (meV) \\
\midrule\noalign{}
\endhead
\bottomrule\noalign{}
\caption{Momentum-mesh convergence of the four representative gap estimates.}\label{tab-nbcp-d-numerical-verification-2}\\
\endlastfoot
12 & Y & PD & 0.0068587113552 \\
12 & Y & Gamma & \(7.29595989041\times10^{-6}\) \\
12 & V & PD & 0.0156410420112 \\
12 & V & Gamma & 0.00203895133135 \\
24 & Y & PD & 0.00686145193025 \\
24 & Y & Gamma & \(7.29600055943\times10^{-6}\) \\
24 & V & PD & 0.0156460286068 \\
24 & V & Gamma & 0.00203898155653 \\
48 & Y & PD & 0.0068621739178 \\
48 & Y & Gamma & \(7.2959923841\times10^{-6}\) \\
48 & V & PD & 0.0156468152768 \\
48 & V & Gamma & 0.00203898166782 \\
\end{longtable}

Local five-point curvature independently checks the angular fits.
Representative results at N=96 and step 0.04 rad:

\begin{longtable}[]{@{}
  >{\raggedright\arraybackslash}p{(\linewidth - 8\tabcolsep) * \real{0.1667}}
  >{\raggedright\arraybackslash}p{(\linewidth - 8\tabcolsep) * \real{0.1667}}
  >{\raggedleft\arraybackslash}p{(\linewidth - 8\tabcolsep) * \real{0.2222}}
  >{\raggedleft\arraybackslash}p{(\linewidth - 8\tabcolsep) * \real{0.2222}}
  >{\raggedleft\arraybackslash}p{(\linewidth - 8\tabcolsep) * \real{0.2222}}@{}}
\toprule\noalign{}
\begin{minipage}[b]{\linewidth}\raggedright
Phase
\end{minipage} & \begin{minipage}[b]{\linewidth}\raggedright
SOC
\end{minipage} & \begin{minipage}[b]{\linewidth}\raggedleft
Gap (meV)
\end{minipage} & \begin{minipage}[b]{\linewidth}\raggedleft
Mesh change
\end{minipage} & \begin{minipage}[b]{\linewidth}\raggedleft
Step change
\end{minipage} \\
\midrule\noalign{}
\endfirsthead
\toprule\noalign{}
\begin{minipage}[b]{\linewidth}\raggedright
Phase
\end{minipage} & \begin{minipage}[b]{\linewidth}\raggedright
SOC
\end{minipage} & \begin{minipage}[b]{\linewidth}\raggedleft
Gap (meV)
\end{minipage} & \begin{minipage}[b]{\linewidth}\raggedleft
Mesh change
\end{minipage} & \begin{minipage}[b]{\linewidth}\raggedleft
Step change
\end{minipage} \\
\midrule\noalign{}
\endhead
\bottomrule\noalign{}
\caption{Local-curvature gap checks and relative mesh and derivative-step changes.}\label{tab-nbcp-d-numerical-verification-3}\\
\endlastfoot
Y & PD & 0.00686213408 & \(1.524\times10^{-5}\) &
\(2.009\times10^{-5}\) \\
Y & Gamma & \(7.29575296\times10^{-6}\) & \(-2.291\times10^{-5}\) &
\(9.395\times10^{-5}\) \\
V & PD & 0.0156467362 & \(8.114\times10^{-6}\) &
\(2.554\times10^{-5}\) \\
V & Gamma & 0.00203897924 & \(-3.107\times10^{-8}\) &
\(1.087\times10^{-6}\) \\
\end{longtable}

Mesh change means \(N=48\to96\) at a fixed step of 0.04 rad; step change
means \(0.04\to0.02\) rad at \(N=96\). The entries are relative changes,
not percentages. Their absolute values are below 0.01\% at these four
representative points. The existing Hamiltonian-pairing regression suite
also passed all eight tests during the calculation session.

For the very weak Y/Gamma energy modulation, making the derivative step
too small amplifies cancellation error; all four step sizes are
preserved in {\small\texttt{\nolinkurl{curvature-validation.json}}}.
These convergence differences do not bound higher-order physical
corrections.

\subsection{Evidence and reproduction}\label{evidence-and-reproduction}

\begin{itemize}
\item
  \href{../../../data-space/verification/260918-y-density-wall/density-wall-check.json}{Classical
  density-wall scan} and
  \href{../../../data-space/verification/260918-y-density-wall/density-wall-validation.json}{refinements},
  reproduced by
  \href{../../../examples/nbcp_y_density_wall.py}{the
  wall diagnostic} and
  \href{../../../examples/nbcp_y_density_wall_validation.py}{the
  refinement script}. These retain full-sphere spin motion, two periodic
  interfaces and explicit phase-boundary conventions. The results are
  classical local minima, not thermal wall free energies.
\item
  \href{../../../data-space/verification/260918-y-nonlinear-gradient/nonlinear-gradient-check.json}{Nonlinear
  smooth-wave results}, reproduced by
  \href{../../../examples/nbcp_y_nonlinear_gradient.py}{the
  constrained classical diagnostic}: 216 configurations, finite-momentum
  harmonic matching, analytic-gradient checks, random starts and a
  fixed-wavelength torus comparison. This tests the classical gradient
  sector within a local Y chart; it does not compute defects or thermal
  coefficients.
\item
  \href{../../../data-space/verification/260918-y-angular-matching/angular-matching-check.json}{Y
  angular matching}: full-circle tensors, principal directions,
  unconstrained harmonics, pure and mixed SOC scans,
  numerical-resolution diagnostics and source hashes. Reproduce with
  \href{../../../examples/nbcp_y_angular_matching.py}{the
  matching diagnostic}.
  \href{../../../data-space/verification/260918-y-angular-matching/independent-check.json}{Independent
  numerical checks}, produced by
  \href{../../../examples/nbcp_y_angular_validation.py}{the
  validation script}, use fresh unaveraged momenta, off-grid tensor
  evaluations and local five-point energy curvatures. These checks do
  not replace quantum or thermal matching or human physics review.
\item
  \href{../../../data-space/verification/260917-y-stability/stability-check.json}{Y
  full-zone stability screen}: mesh and local-refinement minima,
  phase-projector overlaps, adjustable tolerances, direct matrix and
  archived-code comparisons, an intentionally unstable control, runtime
  and input hashes. Reproduce with
  \href{../../../examples/nbcp_y_stability.py}{the
  stability diagnostic}. The calculation tests harmonic stability at one
  field and does not replace a competing-phase or quantum thermal
  calculation.
\item
  \href{../../../data-space/verification/260917-clock-matching/four-case-figure.json}{Four-case
  angular-energy figure}: panel parameters, saved energy differences,
  source hashes and plotting provenance. Reproduce with
  \href{../../../examples/nbcp_clock_matching_figures.py}{the
  four-case plotter}. The original scan arrays and earlier orbit
  illustration remain unchanged.
\item
  \href{../../../data-space/verification/260917-clock-rg/clock-rg-check.json}{Clock
  RG normalization checks}, reproduced by
  \href{../../../examples/nbcp_clock_rg_checks.py}{the
  symbolic check}: shell variance, dipole factors, duality
  normalization, constant-tensor mapping and sixfold exponents. No
  thermal RG trajectory is computed.
\item
  \href{../../../data-space/verification/260917-y-soc-conditions/symbolic-soc-check.json}{Y
  SOC symbolic checks} and
  \href{../../../data-space/verification/260917-y-soc-conditions/soc-conditions-check.json}{numerical
  checks}: first moments, static and dynamical reduction, representative
  coefficients and the bond-direction embedding, now fixed by the
  toolkit convention. Reproduce
  with
  \href{../../../examples/nbcp_y_soc_conditions.wl}{Wolfram
  Language} followed by
  \href{../../../examples/nbcp_y_soc_conditions.py}{Python}.
\item
  \href{../../../data-space/verification/260917-y-stiffness/symbolic-check.json}{Y
  stiffness symbolic checks} and
  \href{../../../data-space/verification/260917-y-stiffness/stiffness-check.json}{numerical
  checks}: the zero-SOC classical baseline and its scope. Reproduce with
  \href{../../../examples/nbcp_y_stiffness.wl}{Wolfram
  Language} followed by
  \href{../../../examples/nbcp_y_stiffness.py}{Python};
  both scripts preserve the existing gap scans.
\item
  \href{../../../data-space/verification/261001-v-mixed-soc/v-mixed-soc-check.json}{V
  mixed-SOC harmonics}: threefold and sixfold amplitudes, phases and
  resolved curvatures near the PD axis, the reflection check and the
  higher harmonics of the saved pure-axis scans. Reproduce with
  \href{../../../examples/nbcp_v_mixed_soc.py}{the mixed-SOC
  diagnostic}. It uses the fixed classical V orbit and the leading
  vacuum energy only.
\item
  \href{../../../data-space/verification/261001-c2t-reflection/c2t-reflection-check.json}{\(C_2'\mathcal T\)
  reflection check}: bond-matrix invariance, the sublattice map of the
  mirror and the vacuum-energy identities \(v(\pi-\phi)=v(\phi)\) and
  \(v(\phi+2\pi/3)=v(\phi)\) on the Y and V orbits for mixed SOC, with
  the reflection about \(\phi=0\) as a control. Reproduce with
  \href{../../../examples/nbcp_c2t_reflection_check.py}{the reflection
  check}.
\item
  \href{../../../data-space/verification/261001-gap-validity/gap-validity-check.json}{Gap
  validity check}: orbit versus straight-line curvature with the
  tadpole term, the uniform six-mode dynamics, stability margins and the
  harmonic finite-temperature curvature (Appendix
  Section~\ref{sec-nbcp-gap-validity}). Reproduce with
  \href{../../../examples/nbcp_gap_validity_check.py}{the validity
  check}.
\item
  \href{../../../data-space/verification/261001-angular-thermal-split/angular-thermal-split.json}{Thermal
  cutoff split} and
  \href{../../../data-space/verification/261002-matching-cutoff/matching-cutoff-check.json}{matching-cutoff
  check}: momentum-resolved harmonic sixfold pinning of Y and V, and its
  comparison with the phase theory built from the relaxed
  \(\rho(\phi)\) (Section~\ref{sec-nbcp-matching-results}). Reproduce
  with \href{../../../examples/nbcp_angular_thermal_split.py}{the split}
  and \href{../../../examples/nbcp_matching_cutoff_check.py}{the
  matching check}.
\item
  \href{../../../data-space/verification/261002-y-vortex-core/vortex-core-check.json}{Vortex-core
  check}: relaxed classical Y vortices of both windings on hexagonal
  clusters up to 6\,937 spins, log coefficient and core energy
  (Section~\ref{sec-nbcp-matching-results}). Reproduce with
  \href{../../../examples/nbcp_y_vortex_core.py}{the vortex script}.
\item
  \href{../../../data-space/verification/261002-one-loop-susceptibility/one-loop-susceptibility.json}{One-loop
  susceptibility}: the \(O(S^0)\) correction to \(\chi_z\) and
  \(m_z\) of Y and V (Appendix Section~\ref{sec-nbcp-gap-validity}).
  Reproduce with
  \href{../../../examples/nbcp_one_loop_susceptibility.py}{the
  susceptibility script}.
\item
  \href{../../../data-space/verification/260916-clock-anisotropy/clock-anisotropy-audit.json}{Clock-anisotropy
  audit JSON}: coefficients at three meshes, fresh matrix checks, input
  and code hashes, and limitations.
\item
  \href{../../../examples/nbcp_clock_anisotropy_audit.py}{Diagnostic
  script}: rerun from the repository root with
  {\small\texttt{\nolinkurl{PYTHONDONTWRITEBYTECODE=1 OPENBLAS_NUM_THREADS=1 python examples/nbcp_clock_anisotropy_audit.py}}}.
\item
  \href{../../../data-space/verification/260912-pseudo-goldstone}{Existing
  scan data}: preserved original energy arrays. New angular and momentum
  refinements are stored separately in the angular-matching directory.
\end{itemize}

\clearpage

\subsection{Reproduction and provenance}\label{sec-reproduction}

\begin{itemize}
\tightlist
\item
  {\small\texttt{\nolinkurl{examples/pseudo_goldstone_comparison.py}}}:
  calculation, checks and original-class replay.
\item
  {\small\texttt{\nolinkurl{examples/pseudo_goldstone_plot.py}}}:
  figures only, regenerated from saved calculation data.
\item
  {\small\texttt{\nolinkurl{examples/pseudo_goldstone_validation.py}}}:
  independent local curvature and \(N=96\) checks.
\item
  {\small\texttt{\nolinkurl{examples/pseudo_goldstone_conjugate_mode.py}}}:
  explicit Y/V partner, constrained relaxation, coordinate invariance
  and weak-pinning dynamics checks.
\item
  {\small\texttt{\nolinkurl{examples/nbcp_research_export.py}}}: PDF
  export from the canonical LaTeX manuscript.
\item
  {\small\texttt{\nolinkurl{scan-N*-P72.json}}}: actual arrays,
  failures, diagnostics and source hashes.
\item
  {\small\texttt{\nolinkurl{curvature-validation.json}}}:
  finite-difference and mesh checks.
\item
  {\small\texttt{\nolinkurl{conjugate-mode-validation-260915.json}}}:
  the 2026-09-15 canonical-partner verification and source hashes.
\item
  {\small\texttt{\nolinkurl{legacy-replay-*.txt}}}: unchanged legacy gap
  class executed with the matched inputs.
\end{itemize}

Run the scan from the repository root:

\begin{Shaded}
\begin{Highlighting}[]
\VariableTok{PYTHONDONTWRITEBYTECODE}\OperatorTok{=}\NormalTok{1 }\VariableTok{OPENBLAS\_NUM\_THREADS}\OperatorTok{=}\NormalTok{1 }\DataTypeTok{\textbackslash{}}
  \ExtensionTok{python}\NormalTok{ examples/pseudo\_goldstone\_comparison.py }\AttributeTok{{-}{-}mesh}\NormalTok{ 48 }\DataTypeTok{\textbackslash{}}
  \AttributeTok{{-}{-}couplings}\NormalTok{ 0 0.0025 0.005 0.0075 0.01 0.0125 0.015 0.0175 0.02}
\end{Highlighting}
\end{Shaded}

To regenerate the PDF preview from the LaTeX source:

\begin{Shaded}
\begin{Highlighting}[]
\ExtensionTok{python}\NormalTok{ examples/nbcp\_research\_export.py}
\end{Highlighting}
\end{Shaded}

The editable source is \texttt{docs/nbcp/main.tex} and its included
chapter, appendix and reference files. Shared formatting and metadata are
separate; \texttt{docs/nbcp/README.md} documents their roles. Original
figures and data remain in \texttt{data-space/verification/}. The XeLaTeX
exporter saves only the PDF and provenance JSON in
\texttt{docs/nbcp/output/}; intermediate build files are temporary.
The former Markdown source and converter are archived.

\subsubsection{Scope and revision
record}\label{scope-and-revision-record}

The numerical comparison was recorded on 2026-09-12, the explicit
canonical-mode derivation and cross-checks on 2026-09-15, and the
research-note layout on 2026-09-16. On 2026-09-16 the detailed gap and
clock notes were consolidated with the main chapters into one Markdown
manuscript. On 2026-09-18 the NBCP manuscript was migrated to chapter-based
LaTeX sources, preserving the equations, results and review status. The
PDF is generated from these TeX sources. The note remains in review and is not an accepted
theory document or a publication-ready result. The production package
and archived scripts were not edited as part of this comparison.

The original gap script defaults to \(J=0.076\) meV and \(h=0.05\) meV.
The comparisons instead supply the same paper parameters and states to
both implementations. An earlier discussion of the internal XXZ orbit
did not invalidate the common-\(z\) orbit used by the archived
SOC-induced Y/V calculation.

The scan field
{\small\texttt{\nolinkurl{old_pair_Berry_coefficient_per_spin}}} records
\(+(S/3)\sum_\alpha\sin\vartheta_\alpha\), opposite to the sign
convention used for \(b_{\mathbf u}\) in the derivation. The gap depends
on its square. The canonical-response evidence and source hashes are
preserved in
{\small\texttt{\nolinkurl{conjugate-mode-validation-260915.json}}}.
